\synctex=1
\documentclass[11pt,aspectratio=169]{beamer}

\usepackage[T1]{fontenc}
\usepackage[french,provide=*]{babel}
\usepackage{amsmath,amssymb,amsfonts}
\usepackage{mathtools}
\usepackage{bm}

% Police serif pour les mathématiques
\usefonttheme[onlymath]{serif}

% Git hash
\usepackage{xstring}
\usepackage{catchfile}
\immediate\write18{git rev-parse HEAD > git.hash}
\CatchFileDef{\HEAD}{git.hash}{\endlinechar=-1}
\newcommand{\gitrevision}{\StrLeft{\HEAD}{7}}

\usepackage{ccicons}

% Pied de page personnalisé
\setbeamertemplate{footline}{
  \hfill\vspace*{1pt}\href{https://creativecommons.org/publicdomain/zero/1.0/deed.fr}{\cczero}\enspace\href{https://github.com/stepan-a/time-series}{\gitrevision}\enspace--\enspace\insertframenumber/\inserttotalframenumber\enspace--\enspace\today\enspace}

% Supprimer la barre de navigation
\setbeamertemplate{navigation symbols}{}

% Commandes personnalisées
\newcommand{\E}{\mathbb{E}}
\newcommand{\Var}{\mathbb{V}}
\newcommand{\Cov}{\mathrm{Cov}}
\newcommand{\Corr}{\mathrm{Corr}}
\newcommand{\R}{\mathbb{R}}
\newcommand{\Z}{\mathbb{Z}}
\newcommand{\N}{\mathbb{N}}

\title{Estimation des modèles ARMA}
\author{}
\date{}
\institute{}

\AtBeginSection[]
{
  \begin{frame}
    \frametitle{Plan}
    \tableofcontents[currentsection]
  \end{frame}
}

\begin{document}

\begin{frame}
\titlepage
\end{frame}

\begin{frame}{Plan}
\tableofcontents
\end{frame}

%==============================================================================
\section{Introduction}
%==============================================================================

\begin{frame}{Le problème de l'estimation}

\begin{itemize}

\item Soit un modèle ARMA($p$, $q$) :
  \[
  Y_t = c + \phi_1 Y_{t-1} + \cdots + \phi_p Y_{t-p} + \varepsilon_t + \theta_1 \varepsilon_{t-1} + \cdots + \theta_q \varepsilon_{t-q}
  \]
  avec $\varepsilon_t$ un bruit blanc.\newline

\item Les chapitres précédents ont montré comment calculer les moments du processus (autocovariances, autocorrélations, prévisions linéaires) en fonction des paramètres.\newline

\item \textbf{Problème :} Comment estimer les paramètres $(c, \phi_1, \ldots, \phi_p, \theta_1, \ldots, \theta_q, \sigma^2)$ à partir d'un échantillon d'observations $(y_1, y_2, \ldots, y_T)$ ?

\end{itemize}

\end{frame}

\begin{frame}{Principe du maximum de vraisemblance}

\begin{itemize}

\item On note $\bm{\theta} = (c, \phi_1, \ldots, \phi_p, \theta_1, \ldots, \theta_q, \sigma^2)'$ le vecteur des paramètres.\newline

\item On suppose que $\varepsilon_t \sim \text{i.i.d.} \; N(0, \sigma^2)$.\newline

\item On calcule la densité jointe de l'échantillon observé :
  \[
  f_{Y_T, Y_{T-1}, \ldots, Y_1}(y_T, y_{T-1}, \ldots, y_1; \bm{\theta})
  \]
  vue comme une fonction de $\bm{\theta}$ pour les données observées.\newline

\item L'\textbf{estimateur du maximum de vraisemblance} (EMV) est la valeur $\hat{\bm{\theta}}$ qui maximise cette fonction, c'est-à-dire la valeur des paramètres pour laquelle l'échantillon observé est le plus probable.

\end{itemize}

\end{frame}

\begin{frame}{Hypothèse de normalité}

\begin{itemize}

\item L'hypothèse de normalité sur $\varepsilon_t$ est forte, mais l'estimation résultante reste pertinente même si elle est violée.\newline

\item Si le vrai processus est non gaussien, les estimateurs obtenus en maximisant la vraisemblance gaussienne restent \textbf{convergents}. On parle alors d'estimateur de \textbf{quasi-maximum de vraisemblance}.\newline

\item En revanche, les erreurs standard calculées sous l'hypothèse de normalité peuvent ne pas être correctes si les données sont non gaussiennes.

\end{itemize}

\end{frame}

\begin{frame}{Décomposition en erreurs de prévision}

\begin{itemize}

\item La densité jointe peut être factorisée en utilisant la règle de Bayes :
  \[
  f_{Y_T, \ldots, Y_1}(y_T, \ldots, y_1; \bm{\theta}) = f_{Y_1}(y_1; \bm{\theta}) \cdot \prod_{t=2}^{T} f_{Y_t | Y_{t-1}}(y_t | y_{t-1}; \bm{\theta})
  \]\newline

\item La \textbf{log-vraisemblance} est donc :
  \[
  \mathcal{L}(\bm{\theta}) = \log f_{Y_1}(y_1; \bm{\theta}) + \sum_{t=2}^{T} \log f_{Y_t | Y_{t-1}}(y_t | y_{t-1}; \bm{\theta})
  \]\newline

\item Cette décomposition est connue sous le nom de \textbf{décomposition en erreurs de prévision} (\emph{prediction-error decomposition}).

\end{itemize}

\end{frame}

%==============================================================================
\section{Vraisemblance d'un processus AR(1) gaussien}
%==============================================================================

\begin{frame}{Le modèle AR(1) gaussien}

\begin{itemize}

\item On considère le processus AR(1) gaussien :
  \[
  Y_t = c + \phi Y_{t-1} + \varepsilon_t
  \]
  avec $\varepsilon_t \sim \text{i.i.d.} \; N(0, \sigma^2)$ et $|\phi| < 1$.\newline

\item Le vecteur de paramètres est $\bm{\theta} = (c, \phi, \sigma^2)'$.\newline

\item L'espérance du processus stationnaire est $\mu = c/(1 - \phi)$.\newline

\item La variance du processus stationnaire est $\sigma^2/(1 - \phi^2)$.

\end{itemize}

\end{frame}

\begin{frame}{Densité de la première observation}

\begin{itemize}

\item Puisque $\varepsilon_t$ est gaussien, $Y_1$ est également gaussien avec :
  \[
  \E[Y_1] = \mu = \frac{c}{1 - \phi}, \quad \Var[Y_1] = \frac{\sigma^2}{1 - \phi^2}
  \]\newline

\item La densité de $Y_1$ est donc :
  \[
  f_{Y_1}(y_1; \bm{\theta}) = \frac{1}{\sqrt{2\pi \sigma^2 / (1 - \phi^2)}} \exp\left[ \frac{-\{y_1 - c/(1 - \phi)\}^2}{2\sigma^2/(1 - \phi^2)} \right]
  \]

\end{itemize}

\end{frame}

\begin{frame}{Densités conditionnelles}

\begin{itemize}

\item Conditionnellement à $Y_{t-1} = y_{t-1}$, on a :
  \[
  Y_t = c + \phi Y_{t-1} + \varepsilon_t \sim N(c + \phi y_{t-1}, \sigma^2)
  \]\newline

\item La densité conditionnelle est :
  \[
  f_{Y_t | Y_{t-1}}(y_t | y_{t-1}; \bm{\theta}) = \frac{1}{\sqrt{2\pi\sigma^2}} \exp\left[ \frac{-(y_t - c - \phi y_{t-1})^2}{2\sigma^2} \right]
  \]\newline

\item Cette expression est valable pour $t = 2, 3, \ldots, T$.

\end{itemize}

\end{frame}

\begin{frame}{Log-vraisemblance exacte}

\begin{itemize}

\item La log-vraisemblance exacte de l'AR(1) gaussien est :
  \begin{align*}
  \mathcal{L}(\bm{\theta}) &= -\frac{1}{2}\log(2\pi) - \frac{1}{2}\log\!\left(\frac{\sigma^2}{1 - \phi^2}\right) - \frac{\{y_1 - c/(1 - \phi)\}^2}{2\sigma^2/(1 - \phi^2)} \\
  &\quad - \frac{T-1}{2}\log(2\pi) - \frac{T-1}{2}\log(\sigma^2) \\
  &\quad - \sum_{t=2}^{T} \frac{(y_t - c - \phi y_{t-1})^2}{2\sigma^2}
  \end{align*}

\item On regroupe les termes pour obtenir :
  \begin{align*}
  \mathcal{L}(\bm{\theta}) &= -\frac{T}{2}\log(2\pi) - \frac{T}{2}\log(\sigma^2) + \frac{1}{2}\log(1 - \phi^2) \\
  &\quad - \frac{(1 - \phi^2)}{2\sigma^2}\left(y_1 - \frac{c}{1 - \phi}\right)^2 - \sum_{t=2}^{T} \frac{(y_t - c - \phi y_{t-1})^2}{2\sigma^2}
  \end{align*}

\end{itemize}

\end{frame}

\begin{frame}{Log-vraisemblance conditionnelle}

\begin{itemize}

\item Une alternative consiste à conditionner sur la première observation $y_1$ et à maximiser :
  \[
  \mathcal{L}_c(\bm{\theta}) = -\frac{T-1}{2}\log(2\pi) - \frac{T-1}{2}\log(\sigma^2) - \sum_{t=2}^{T} \frac{(y_t - c - \phi y_{t-1})^2}{2\sigma^2}
  \]\newline

\item La maximisation par rapport à $c$ et $\phi$ revient à minimiser :
  \[
  \sum_{t=2}^{T} (y_t - c - \phi y_{t-1})^2
  \]\newline

\item C'est un problème de \textbf{moindres carrés ordinaires} (MCO) : régression de $y_t$ sur une constante et $y_{t-1}$.

\end{itemize}

\end{frame}

\begin{frame}{Estimateurs conditionnels du maximum de vraisemblance}

\begin{itemize}

\item Les estimateurs conditionnels de $c$ et $\phi$ sont donnés par la formule des MCO :
  \[
  \begin{bmatrix} \hat{c} \\ \hat{\phi} \end{bmatrix} = \begin{bmatrix} T-1 & \sum y_{t-1} \\ \sum y_{t-1} & \sum y_{t-1}^2 \end{bmatrix}^{-1} \begin{bmatrix} \sum y_t \\ \sum y_{t-1} y_t \end{bmatrix}
  \]
  où les sommes portent sur $t = 2, 3, \ldots, T$.\newline

\item L'estimateur conditionnel de $\sigma^2$ est la variance résiduelle de la régression :
  \[
  \hat{\sigma}^2 = \frac{1}{T-1} \sum_{t=2}^{T} (y_t - \hat{c} - \hat{\phi} y_{t-1})^2
  \]\newline

\item Pour $T$ grand, les estimateurs exact et conditionnel convergent vers la même distribution asymptotique (si $|\phi| < 1$).

\end{itemize}

\end{frame}

\begin{frame}{Vraisemblance exacte et conditionnelle : comparaison}

\begin{itemize}

\item La vraisemblance exacte nécessite le terme supplémentaire $\frac{1}{2}\log(1 - \phi^2) - \frac{(1-\phi^2)}{2\sigma^2}(y_1 - \mu)^2$ lié à la première observation.\newline

\item La vraisemblance conditionnelle ignore ce terme et traite $y_1$ comme déterministe.\newline

\item Si $T$ est grand, la contribution de la première observation est négligeable.\newline

\item Si $|\phi| < 1$, les deux approches donnent des estimateurs \textbf{convergents}. Si $|\phi| > 1$, la vraisemblance conditionnelle ne fournit pas d'estimateurs convergents car la densité de $Y_1$ n'est pas correctement spécifiée.

\end{itemize}

\end{frame}

\begin{frame}{Expression vectorielle de la vraisemblance (1/6)}

\begin{itemize}

\item On peut dériver la vraisemblance d'une manière alternative en considérant le vecteur des $T$ observations :
  \[
  \mathbf{y} = (y_1, y_2, \ldots, y_T)'
  \]\newline

\item Ce vecteur peut être vu comme une unique réalisation d'un vecteur gaussien de dimension $T$ :
  \[
  \mathbf{y} \sim N(\bm{\mu}, \bm{\Omega})
  \]\newline

\item Le vecteur d'espérances est $\bm{\mu} = (\mu, \mu, \ldots, \mu)'$ avec $\mu = c/(1 - \phi)$.\newline

\item La matrice de variance-covariance $\bm{\Omega}$ est une matrice $(T \times T)$ dont l'élément $(i, j)$ est l'autocovariance $\gamma(|i-j|)$.

\end{itemize}

\end{frame}

\begin{frame}{Expression vectorielle de la vraisemblance (2/6)}

\begin{itemize}

\item Pour l'AR(1), l'autocovariance est $\gamma(h) = \sigma^2 \phi^{|h|}/(1-\phi^2)$, donc :
  \[
  \bm{\Omega} = \sigma^2 \mathbf{V}
  \]
  où
  \[
  \mathbf{V} = \frac{1}{1-\phi^2} \begin{bmatrix} 1 & \phi & \phi^2 & \cdots & \phi^{T-1} \\ \phi & 1 & \phi & \cdots & \phi^{T-2} \\ \phi^2 & \phi & 1 & \cdots & \phi^{T-3} \\ \vdots & \vdots & \vdots & \ddots & \vdots \\ \phi^{T-1} & \phi^{T-2} & \phi^{T-3} & \cdots & 1 \end{bmatrix}
  \]\newline

\item $\mathbf{V}$ est une matrice de Toeplitz symétrique définie positive.

\end{itemize}

\end{frame}

\begin{frame}{Expression vectorielle de la vraisemblance (3/6)}

\begin{itemize}

\item La densité du vecteur gaussien $\mathbf{y}$ s'écrit :
  \[
  f_{\mathbf{Y}}(\mathbf{y}; \bm{\theta}) = (2\pi)^{-T/2} |\bm{\Omega}|^{-1/2} \exp\!\left[-\tfrac{1}{2}(\mathbf{y} - \bm{\mu})'\bm{\Omega}^{-1}(\mathbf{y} - \bm{\mu})\right]
  \]\newline

\item La log-vraisemblance est donc :
  \[
  \boxed{\mathcal{L}(\bm{\theta}) = -\frac{T}{2}\log(2\pi) + \frac{1}{2}\log|\bm{\Omega}^{-1}| - \frac{1}{2}(\mathbf{y} - \bm{\mu})'\bm{\Omega}^{-1}(\mathbf{y} - \bm{\mu})}
  \]\newline

\item Cette expression fait intervenir l'inverse et le déterminant de la matrice $\bm{\Omega}$ de taille $(T \times T)$. En apparence, le calcul est coûteux, mais la structure de $\bm{\Omega}$ permet de le simplifier.

\end{itemize}

\end{frame}

\begin{frame}{Expression vectorielle de la vraisemblance (4/6)}

\begin{itemize}

\item On introduit la matrice triangulaire inférieure $\mathbf{L}$ de taille $(T \times T)$ :
  \[
  \mathbf{L} = \begin{bmatrix} \sqrt{1-\phi^2} & 0 & 0 & \cdots & 0 & 0 \\ -\phi & 1 & 0 & \cdots & 0 & 0 \\ 0 & -\phi & 1 & \cdots & 0 & 0 \\ \vdots & \vdots & \vdots & \ddots & \vdots & \vdots \\ 0 & 0 & 0 & \cdots & -\phi & 1 \end{bmatrix}
  \]\newline

\item On peut montrer que :
  \[
  \mathbf{L}'\mathbf{L} = \mathbf{V}^{-1}
  \]
  et donc $\bm{\Omega}^{-1} = \sigma^{-2}\mathbf{V}^{-1} = \sigma^{-2}\mathbf{L}'\mathbf{L}$.

\end{itemize}

\end{frame}

\begin{frame}{Expression vectorielle de la vraisemblance (5/6)}

\begin{itemize}

\item Puisque $\mathbf{L}$ est triangulaire, son déterminant est le produit des éléments diagonaux :
  \[
  |\mathbf{L}| = \sqrt{1-\phi^2} \cdot \underbrace{1 \cdot 1 \cdots 1}_{T-1} = \sqrt{1-\phi^2}
  \]\newline

\item Donc $|\mathbf{L}'\mathbf{L}| = |\mathbf{V}^{-1}| = 1 - \phi^2$, et :
  \[
  \frac{1}{2}\log|\bm{\Omega}^{-1}| = \frac{1}{2}\log(\sigma^{-2T} \cdot |\mathbf{V}^{-1}|) = -\frac{T}{2}\log(\sigma^2) + \frac{1}{2}\log(1 - \phi^2)
  \]\newline

\item On définit le vecteur transformé $\tilde{\mathbf{y}} = \mathbf{L}(\mathbf{y} - \bm{\mu})$, de sorte que :
  \[
  (\mathbf{y} - \bm{\mu})'\bm{\Omega}^{-1}(\mathbf{y} - \bm{\mu}) = \frac{1}{\sigma^2}\tilde{\mathbf{y}}'\tilde{\mathbf{y}}
  \]

\end{itemize}

\end{frame}

\begin{frame}{Expression vectorielle de la vraisemblance (6/6)}

\begin{itemize}

\item Les composantes de $\tilde{\mathbf{y}} = \mathbf{L}(\mathbf{y} - \bm{\mu})$ sont :
  \[
  \tilde{y}_1 = \sqrt{1-\phi^2}\,(y_1 - \mu), \quad \tilde{y}_t = (y_t - \mu) - \phi(y_{t-1} - \mu) \;\; \text{pour } t \geq 2
  \]\newline

\item En substituant $\mu = c/(1-\phi)$ :
  \[
  \tilde{y}_t = y_t - c - \phi y_{t-1} \quad \text{pour } t \geq 2
  \]
  Ce sont les \textbf{erreurs de prévision} !\newline

\item La log-vraisemblance s'écrit alors :
  \[
  \boxed{\mathcal{L}(\bm{\theta}) = -\frac{T}{2}\log(2\pi) - \frac{T}{2}\log(\sigma^2) + \frac{1}{2}\log(1-\phi^2) - \frac{1}{2\sigma^2}\sum_{t=1}^{T} \tilde{y}_t^2}
  \]\newline

\item On retrouve exactement l'expression de la log-vraisemblance exacte obtenue par la décomposition en erreurs de prévision.

\end{itemize}

\end{frame}

%==============================================================================
\section{Vraisemblance d'un processus AR($p$) gaussien}
%==============================================================================

\begin{frame}{Le modèle AR($p$) gaussien}

\begin{itemize}

\item On considère le processus AR($p$) gaussien :
  \[
  Y_t = c + \phi_1 Y_{t-1} + \phi_2 Y_{t-2} + \cdots + \phi_p Y_{t-p} + \varepsilon_t
  \]
  avec $\varepsilon_t \sim \text{i.i.d.}\; N(0, \sigma^2)$.\newline

\item Le vecteur de paramètres est $\bm{\theta} = (c, \phi_1, \phi_2, \ldots, \phi_p, \sigma^2)'$.\newline

\item L'espérance du processus stationnaire est :
  \[
  \mu = \frac{c}{1 - \phi_1 - \phi_2 - \cdots - \phi_p}
  \]

\end{itemize}

\end{frame}

\begin{frame}{Construction de la vraisemblance}

\begin{itemize}

\item Les $p$ premières observations $(y_1, \ldots, y_p)$ suivent une loi gaussienne multivariée $N(\bm{\mu}_p, \sigma^2 \mathbf{V}_p)$.\newline

\item Pour $t > p$, conditionnellement aux $p$ observations précédentes :
  \[
  Y_t | Y_{t-1}, \ldots, Y_{t-p} \sim N(c + \phi_1 y_{t-1} + \cdots + \phi_p y_{t-p}, \; \sigma^2)
  \]\newline

\item La log-vraisemblance exacte est :
  \begin{align*}
  \mathcal{L}(\bm{\theta}) &= -\frac{T}{2}\log(2\pi) - \frac{T}{2}\log(\sigma^2) + \frac{1}{2}\log|\mathbf{V}_p^{-1}| \\
  &\quad - \frac{1}{2\sigma^2}(\mathbf{y}_p - \bm{\mu}_p)'\mathbf{V}_p^{-1}(\mathbf{y}_p - \bm{\mu}_p) \\
  &\quad - \sum_{t=p+1}^{T} \frac{(y_t - c - \phi_1 y_{t-1} - \cdots - \phi_p y_{t-p})^2}{2\sigma^2}
  \end{align*}

\end{itemize}

\end{frame}

\begin{frame}{Estimateurs conditionnels pour l'AR($p$)}

\begin{itemize}

\item La log-vraisemblance conditionnelle (sur les $p$ premières observations) est :
  \begin{align*}
  \mathcal{L}_c(\bm{\theta}) &= -\frac{T-p}{2}\log(2\pi) - \frac{T-p}{2}\log(\sigma^2) \\
  &\quad - \sum_{t=p+1}^{T} \frac{(y_t - c - \phi_1 y_{t-1} - \cdots - \phi_p y_{t-p})^2}{2\sigma^2}
  \end{align*}\newline

\item La maximisation revient à minimiser :
  \[
  \sum_{t=p+1}^{T} (y_t - c - \phi_1 y_{t-1} - \cdots - \phi_p y_{t-p})^2
  \]\newline

\item C'est la somme des carrés des résidus d'une régression MCO de $y_t$ sur une constante et ses $p$ valeurs retardées.

\end{itemize}

\end{frame}

\begin{frame}{Estimateur de la variance des innovations}

\begin{itemize}

\item L'estimateur conditionnel de $\sigma^2$ est le résidu quadratique moyen :
  \[
  \hat{\sigma}^2 = \frac{1}{T-p} \sum_{t=p+1}^{T} (y_t - \hat{c} - \hat{\phi}_1 y_{t-1} - \hat{\phi}_2 y_{t-2} - \cdots - \hat{\phi}_p y_{t-p})^2
  \]\newline

\item Les estimateurs exact et conditionnel ont la même distribution asymptotique.\newline

\item \textbf{Résultat important :} Pour un processus AR($p$), les estimateurs conditionnels du maximum de vraisemblance sont identiques aux estimateurs des MCO. L'estimation d'un AR($p$) est donc particulièrement simple.

\end{itemize}

\end{frame}

%==============================================================================
\section{Estimation par quasi-maximum de vraisemblance}
%==============================================================================

\begin{frame}{Processus non gaussiens}

\begin{itemize}

\item Que se passe-t-il si le processus n'est pas gaussien ?\newline

\item La régression MCO de $y_t$ sur une constante et ses $p$ retards fournit une estimation convergente de la \textbf{projection linéaire} :
  \[
  \hat{\E}(Y_t | Y_{t-1}, Y_{t-2}, \ldots, Y_{t-p})
  \]
  pourvu que le processus soit ergodique pour les moments d'ordre 2.\newline

\item Cette régression MCO maximise aussi la \textbf{log-vraisemblance gaussienne conditionnelle}. Même si le processus n'est pas gaussien, maximiser cette fonction fournit des estimateurs convergents.

\end{itemize}

\end{frame}

\begin{frame}{Estimateur de quasi-maximum de vraisemblance}

\begin{itemize}

\item Un estimateur obtenu en maximisant une vraisemblance mal spécifiée (par exemple, gaussienne alors que les données ne le sont pas) est appelé estimateur de \textbf{quasi-maximum de vraisemblance} (QMLE).\newline

\item \textbf{Propriété :} Le QMLE fournit des estimateurs \textbf{convergents} des paramètres $(\phi_1, \ldots, \phi_p)$.\newline

\item Cependant, les erreurs standard calculées sous l'hypothèse gaussienne peuvent ne pas être correctes. Il faut utiliser un estimateur robuste de la matrice de variance-covariance.\newline

\item En pratique, si les données sont non gaussiennes, une transformation préalable (par exemple, le logarithme) peut rapprocher la distribution de la normalité.

\end{itemize}

\end{frame}

\begin{frame}{Transformations Box-Cox}

\begin{itemize}

\item Pour une variable positive $Y_t$, Box et Cox (1964) proposent la famille de transformations :
  \[
  Y_t^{(\lambda)} = \begin{cases} \dfrac{Y_t^{\lambda} - 1}{\lambda} & \text{si } \lambda \neq 0 \\[0.3cm] \log Y_t & \text{si } \lambda = 0 \end{cases}
  \]\newline

\item On choisit $\lambda$ de sorte que $Y_t^{(\lambda)}$ soit bien approximé par un processus ARMA gaussien.\newline

\item En pratique, pour les séries économiques qui croissent dans le temps (PIB, prix), on utilise souvent :
  \[
  y_t = \log X_t - \log X_{t-1}
  \]
  c'est-à-dire le taux de croissance en log.

\end{itemize}

\end{frame}

%==============================================================================
\section{Vraisemblance d'un processus MA(1) gaussien}
%==============================================================================

\begin{frame}{Le modèle MA(1) gaussien}

\begin{itemize}

\item On considère le processus MA(1) gaussien :
  \[
  Y_t = \mu + \varepsilon_t + \theta \varepsilon_{t-1}
  \]
  avec $\varepsilon_t \sim \text{i.i.d.}\; N(0, \sigma^2)$.\newline

\item Le vecteur de paramètres est $\bm{\theta} = (\mu, \theta, \sigma^2)'$.\newline

\item Contrairement à l'AR, la vraisemblance du MA ne se réduit \textbf{pas} à un problème de moindres carrés.\newline

\item Deux approches : la vraisemblance \textbf{conditionnelle} (plus simple) et la vraisemblance \textbf{exacte} (plus précise en petit échantillon).

\end{itemize}

\end{frame}

\begin{frame}{Vraisemblance conditionnelle du MA(1) (1/2)}

\begin{itemize}

\item On conditionne sur la valeur initiale $\varepsilon_0 = 0$. Sachant $\varepsilon_{t-1}$, la densité de $Y_t$ est :
  \[
  f_{Y_t | \varepsilon_{t-1}}(y_t | \varepsilon_{t-1}; \bm{\theta}) = \frac{1}{\sqrt{2\pi\sigma^2}} \exp\left[\frac{-(y_t - \mu - \theta \varepsilon_{t-1})^2}{2\sigma^2}\right]
  \]\newline

\item Sachant $\varepsilon_0 = 0$, on déduit $\varepsilon_1$ de l'observation $y_1$ :
  \[
  \varepsilon_1 = y_1 - \mu
  \]\newline

\item Plus généralement, les innovations sont calculées récursivement :
  \[
  \varepsilon_t = y_t - \mu - \theta \varepsilon_{t-1}
  \]
  pour $t = 1, 2, \ldots, T$, en partant de $\varepsilon_0 = 0$.

\end{itemize}

\end{frame}

\begin{frame}{Vraisemblance conditionnelle du MA(1) (2/2)}

\begin{itemize}

\item La log-vraisemblance conditionnelle est :
  \[
  \mathcal{L}_c(\bm{\theta}) = -\frac{T}{2}\log(2\pi) - \frac{T}{2}\log(\sigma^2) - \sum_{t=1}^{T} \frac{\varepsilon_t^2}{2\sigma^2}
  \]
  où $\varepsilon_t = y_t - \mu - \theta \varepsilon_{t-1}$ avec $\varepsilon_0 = 0$.\newline

\item La log-vraisemblance est une fonction \textbf{non linéaire} de $\mu$ et $\theta$ : pas de solution analytique explicite.\newline

\item La maximisation requiert une \textbf{optimisation numérique}.\newline

\item \textbf{Condition d'inversibilité :} L'approximation conditionnelle est valide si $|\theta| < 1$. Si l'optimisation conduit à $|\hat{\theta}| > 1$, il faut recommencer avec $\hat{\theta}^{-1}$.

\end{itemize}

\end{frame}

\begin{frame}{Effet de la condition initiale}

\begin{itemize}

\item En développant la récurrence $\varepsilon_t = y_t - \mu - \theta \varepsilon_{t-1}$ :
  \begin{align*}
  \varepsilon_t &= (y_t - \mu) - \theta(y_{t-1} - \mu) + \theta^2(y_{t-2} - \mu) - \cdots \\
  &\quad + (-1)^{t-1} \theta^{t-1}(y_1 - \mu) + (-1)^t \theta^t \varepsilon_0
  \end{align*}\newline

\item Si $|\theta| < 1$, l'effet de $\varepsilon_0 = 0$ décroît géométriquement : $(-1)^t \theta^t \varepsilon_0 \to 0$.\newline

\item Pour $T$ suffisamment grand, l'approximation $\varepsilon_0 = 0$ est inoffensive.\newline

\item Si $|\theta|$ est proche de 1, les premières innovations $\varepsilon_t$ sont mal estimées, ce qui peut affecter l'estimation en petit échantillon.

\end{itemize}

\end{frame}

\begin{frame}{Vraisemblance exacte du MA(1) (1/6)}

\begin{itemize}

\item La vraisemblance exacte traite les $T$ observations comme un vecteur gaussien $\mathbf{y} = (y_1, \ldots, y_T)'$ de loi $N(\bm{\mu}, \bm{\Omega})$.\newline

\item La matrice de variance-covariance $\bm{\Omega}$ a une structure \textbf{tridiagonale} :
  \[
  \bm{\Omega} = \sigma^2 \begin{bmatrix} 1+\theta^2 & \theta & 0 & \cdots & 0 \\ \theta & 1+\theta^2 & \theta & \cdots & 0 \\ 0 & \theta & 1+\theta^2 & \cdots & 0 \\ \vdots & \vdots & \vdots & \ddots & \vdots \\ 0 & 0 & 0 & \cdots & 1+\theta^2 \end{bmatrix}
  \]\newline

\item La log-vraisemblance exacte est :
  \[
  \mathcal{L}(\bm{\theta}) = -\frac{T}{2}\log(2\pi) - \frac{1}{2}\log|\bm{\Omega}| - \frac{1}{2}(\mathbf{y} - \bm{\mu})'\bm{\Omega}^{-1}(\mathbf{y} - \bm{\mu})
  \]

\end{itemize}

\end{frame}

\begin{frame}{Vraisemblance exacte du MA(1) (2/6)}

\begin{itemize}

\item On cherche la \textbf{factorisation triangulaire} $\bm{\Omega} = \mathbf{A}\mathbf{D}\mathbf{A}'$ où :
  \begin{itemize}
    \item $\mathbf{A}$ est triangulaire inférieure avec des 1 sur la diagonale,
    \item $\mathbf{D}$ est diagonale avec des éléments strictement positifs.
  \end{itemize}

\item Puisque $\bm{\Omega}$ est tridiagonale, $\mathbf{A}$ est \textbf{bidiagonale} : seuls les éléments diagonaux et sous-diagonaux sont non nuls.\newline

\item On note $S_t = 1 + \theta^2 + \theta^4 + \cdots + \theta^{2(t-1)}$ la somme géométrique. On a $S_1 = 1$ et :
  \[
  S_t = \frac{1 - \theta^{2t}}{1 - \theta^2} \quad \text{si } \theta^2 \neq 1
  \]

\end{itemize}

\end{frame}

\begin{frame}{Vraisemblance exacte du MA(1) (3/6)}

\begin{itemize}

\item On obtient $\mathbf{A}$ et $\mathbf{D}$ par \textbf{élimination de Gauss} sur $\bm{\Omega}/\sigma^2$.\newline

\item \textbf{Étape 1 :} Le pivot est $d_1 = (1 + \theta^2) = S_2/S_1$. On élimine le terme sous-diagonal $\theta$ en soustrayant $\frac{\theta}{1+\theta^2}$ fois la première ligne. Cela donne $a_{21} = \frac{\theta}{1+\theta^2} = \frac{\theta S_1}{S_2}$.\newline

\item \textbf{Étape 2 :} Le nouveau pivot est :
  \[
  d_2 = (1+\theta^2) - \frac{\theta^2}{1+\theta^2} = \frac{(1+\theta^2)^2 - \theta^2}{1+\theta^2} = \frac{1+\theta^2+\theta^4}{1+\theta^2} = \frac{S_3}{S_2}
  \]
  Le multiplicateur est $a_{32} = \frac{\theta}{d_2} = \frac{\theta S_2}{S_3}$.

\end{itemize}

\end{frame}

\begin{frame}{Vraisemblance exacte du MA(1) (4/6)}

\begin{itemize}

\item En poursuivant, on montre par récurrence que pour $t = 1, \ldots, T$ :
  \[
  \boxed{d_t = \frac{S_{t+1}}{S_t}} \quad \text{et} \quad \boxed{a_{t+1,t} = \frac{\theta \, S_t}{S_{t+1}}}
  \]\newline

\item Explicitement, les matrices sont (en factorisant $\sigma^2$ dans $\mathbf{D}$) :
  \[
  \mathbf{A} = \begin{bmatrix} 1 & 0 & 0 & \cdots & 0 \\[3pt] \frac{\theta S_1}{S_2} & 1 & 0 & \cdots & 0 \\[3pt] 0 & \frac{\theta S_2}{S_3} & 1 & \cdots & 0 \\[3pt] \vdots & \vdots & \vdots & \ddots & \vdots \\[3pt] 0 & 0 & 0 & \cdots & 1 \end{bmatrix}, \quad \mathbf{D} = \sigma^2 \begin{bmatrix} \frac{S_2}{S_1} & & \\[3pt] & \frac{S_3}{S_2} & \\[3pt] & & \ddots \\[3pt] & & & \frac{S_{T+1}}{S_T} \end{bmatrix}
  \]

\end{itemize}

\end{frame}

\begin{frame}{Vraisemblance exacte du MA(1) (5/6)}

\begin{itemize}

\item Le vecteur transformé $\tilde{\mathbf{y}} = \mathbf{A}^{-1}(\mathbf{y} - \bm{\mu})$ a pour composantes :
  \[
  \tilde{y}_1 = y_1 - \mu, \qquad \tilde{y}_t = (y_t - \mu) - \frac{\theta \, S_{t-1}}{S_t}\,\tilde{y}_{t-1} \quad \text{pour } t \geq 2
  \]\newline

\item Puisque $|\mathbf{A}| = 1$ (triangulaire avec des 1 sur la diagonale) et $\bm{\Omega} = \mathbf{A}\mathbf{D}\mathbf{A}'$ :
  \[
  |\bm{\Omega}| = |\mathbf{D}| = \sigma^{2T} \prod_{t=1}^{T} \frac{S_{t+1}}{S_t} = \sigma^{2T}\, \frac{S_{T+1}}{S_1} = \sigma^{2T}\, S_{T+1}
  \]
  car le produit est télescopique.\newline

\item De plus :
  \[
  (\mathbf{y} - \bm{\mu})'\bm{\Omega}^{-1}(\mathbf{y} - \bm{\mu}) = \tilde{\mathbf{y}}'\mathbf{D}^{-1}\tilde{\mathbf{y}} = \sum_{t=1}^{T} \frac{\tilde{y}_t^2}{d_t}
  \]

\end{itemize}

\end{frame}

\begin{frame}{Vraisemblance exacte du MA(1) (6/6)}

\begin{itemize}

\item La log-vraisemblance exacte s'écrit alors :
  \[
  \boxed{\mathcal{L}(\bm{\theta}) = -\frac{T}{2}\log(2\pi) - \frac{1}{2}\sum_{t=1}^{T}\log(d_t) - \frac{1}{2}\sum_{t=1}^{T}\frac{\tilde{y}_t^2}{d_t}}
  \]
  avec $d_t = \sigma^2 S_{t+1}/S_t$ et $S_t = 1 + \theta^2 + \cdots + \theta^{2(t-1)}$.\newline

\item Cette expression est valide pour \textbf{toute} valeur de $\theta$ (pas seulement $|\theta| < 1$).\newline

\item On montre que si $\theta = \hat{\theta}$ maximise cette expression, alors $\hat{\theta}^{-1}$ donne la même valeur. On retient la solution inversible : $|\hat{\theta}| < 1$.\newline

\item Pour la vraisemblance conditionnelle, cette propriété n'est pas garantie.

\end{itemize}

\end{frame}

%==============================================================================
\section{Vraisemblance d'un processus MA($q$) gaussien}
%==============================================================================

\begin{frame}{Le modèle MA($q$) gaussien}

\begin{itemize}

\item On considère le processus MA($q$) gaussien :
  \[
  Y_t = \mu + \varepsilon_t + \theta_1 \varepsilon_{t-1} + \theta_2 \varepsilon_{t-2} + \cdots + \theta_q \varepsilon_{t-q}
  \]
  avec $\varepsilon_t \sim \text{i.i.d.}\; N(0, \sigma^2)$.\newline

\item Le vecteur de paramètres est $\bm{\theta} = (\mu, \theta_1, \ldots, \theta_q, \sigma^2)'$.\newline

\item La matrice de variance-covariance $\bm{\Omega}$ est une matrice bande de largeur $q$ :
  les autocovariances $\gamma_k$ sont nulles pour $k > q$.

\end{itemize}

\end{frame}

\begin{frame}{Vraisemblance conditionnelle du MA($q$)}

\begin{itemize}

\item On pose $\varepsilon_0 = \varepsilon_{-1} = \cdots = \varepsilon_{-q+1} = 0$.\newline

\item Les innovations sont calculées par récurrence :
  \[
  \varepsilon_t = y_t - \mu - \theta_1 \varepsilon_{t-1} - \theta_2 \varepsilon_{t-2} - \cdots - \theta_q \varepsilon_{t-q}
  \]
  pour $t = 1, 2, \ldots, T$.\newline

\item La log-vraisemblance conditionnelle est :
  \[
  \mathcal{L}_c(\bm{\theta}) = -\frac{T}{2}\log(2\pi) - \frac{T}{2}\log(\sigma^2) - \sum_{t=1}^{T}\frac{\varepsilon_t^2}{2\sigma^2}
  \]\newline

\item Cette approximation est valide si toutes les racines de $1 + \theta_1 z + \cdots + \theta_q z^q = 0$ sont de module $> 1$ (inversibilité).

\end{itemize}

\end{frame}

\begin{frame}{Vraisemblance exacte du MA($q$)}

\begin{itemize}

\item La structure bande de $\bm{\Omega}$ permet d'utiliser la factorisation triangulaire $\bm{\Omega} = \mathbf{A}\mathbf{D}\mathbf{A}'$.\newline

\item $\mathbf{A}$ est une matrice triangulaire inférieure bande : $a_{ij} = 0$ pour $i > q + j$.\newline

\item Les éléments de $\tilde{\mathbf{y}} = \mathbf{A}^{-1}(\mathbf{y} - \bm{\mu})$ se calculent récursivement par résolution d'un système triangulaire.\newline

\item La log-vraisemblance exacte est :
  \[
  \mathcal{L}(\bm{\theta}) = -\frac{T}{2}\log(2\pi) - \frac{1}{2}\sum_{t=1}^{T}\log(d_{tt}) - \frac{1}{2}\sum_{t=1}^{T}\frac{\tilde{y}_t^2}{d_{tt}}
  \]\newline

\item Contrairement à la vraisemblance conditionnelle, l'expression exacte est valide pour toute valeur de $(\theta_1, \ldots, \theta_q)$.

\end{itemize}

\end{frame}

%==============================================================================
\section{Vraisemblance d'un processus ARMA($p$, $q$) gaussien}
%==============================================================================

\begin{frame}{Le modèle ARMA($p$, $q$) gaussien}

\begin{itemize}

\item Le modèle ARMA($p$, $q$) gaussien s'écrit :
  \[
  Y_t = c + \phi_1 Y_{t-1} + \cdots + \phi_p Y_{t-p} + \varepsilon_t + \theta_1 \varepsilon_{t-1} + \cdots + \theta_q \varepsilon_{t-q}
  \]
  avec $\varepsilon_t \sim \text{i.i.d.}\; N(0, \sigma^2)$.\newline

\item Le vecteur de paramètres est $\bm{\theta} = (c, \phi_1, \ldots, \phi_p, \theta_1, \ldots, \theta_q, \sigma^2)'$.\newline

\item Le nombre de paramètres est $p + q + 2$.

\end{itemize}

\end{frame}

\begin{frame}{Vraisemblance conditionnelle du ARMA($p$, $q$)}

\begin{itemize}

\item On fixe les conditions initiales :
  \begin{itemize}
  \item $\mathbf{y}_0 = (y_0, y_{-1}, \ldots, y_{-p+1})'$ aux valeurs observées ou à l'espérance,
  \item $\bm{\varepsilon}_0 = (\varepsilon_0, \varepsilon_{-1}, \ldots, \varepsilon_{-q+1})' = \mathbf{0}$.
  \end{itemize}

\item Les innovations sont calculées récursivement :
  \[
  \varepsilon_t = y_t - c - \phi_1 y_{t-1} - \cdots - \phi_p y_{t-p} - \theta_1 \varepsilon_{t-1} - \cdots - \theta_q \varepsilon_{t-q}
  \]
  pour $t = 1, 2, \ldots, T$.\newline

\item La log-vraisemblance conditionnelle est :
  \[
  \mathcal{L}_c(\bm{\theta}) = -\frac{T}{2}\log(2\pi) - \frac{T}{2}\log(\sigma^2) - \sum_{t=1}^{T} \frac{\varepsilon_t^2}{2\sigma^2}
  \]

\end{itemize}

\end{frame}

\begin{frame}{Conditions initiales alternatives}

\begin{itemize}

\item \textbf{Approche Box-Jenkins :} On fixe $y_s = c/(1 - \phi_1 - \cdots - \phi_p)$ pour $s \leq 0$ et $\varepsilon_s = 0$ pour $s \leq 0$.\newline

\item \textbf{Approche alternative :} On fixe les $\varepsilon$ à zéro et les $y$ à leurs valeurs observées. On commence l'itération à $t = p + 1$ :
  \[
  \varepsilon_p = \varepsilon_{p-1} = \cdots = \varepsilon_{p-q+1} = 0
  \]
  La log-vraisemblance conditionnelle est alors :
  \[
  \mathcal{L}_c(\bm{\theta}) = -\frac{T-p}{2}\log(2\pi) - \frac{T-p}{2}\log(\sigma^2) - \sum_{t=p+1}^{T}\frac{\varepsilon_t^2}{2\sigma^2}
  \]\newline

\item Ces approximations sont valides si toutes les racines de $\phi(z) = 0$ et $\theta(z) = 0$ sont de module $> 1$.

\end{itemize}

\end{frame}

\begin{frame}{Vraisemblance exacte}

\begin{itemize}

\item L'approche la plus rigoureuse utilise le \textbf{filtre de Kalman} pour calculer la vraisemblance exacte.\newline

\item Le filtre de Kalman fournit, de manière récursive, les prévisions optimales et les erreurs de prévision, ce qui permet de construire la vraisemblance par décomposition en erreurs de prévision.\newline

\item Alternativement, on peut utiliser la \textbf{factorisation triangulaire} de la matrice de variance-covariance $\bm{\Omega}$, comme dans le cas MA.\newline

\item Les deux approches donnent le même résultat mais diffèrent en termes d'implémentation informatique.

\end{itemize}

\end{frame}

%==============================================================================
\section{Optimisation numérique}
%==============================================================================

\begin{frame}{Le problème d'optimisation}

\begin{itemize}

\item La log-vraisemblance est une fonction non linéaire de $\bm{\theta}$. Il faut trouver :
  \[
  \hat{\bm{\theta}} = \arg\max_{\bm{\theta}} \mathcal{L}(\bm{\theta})
  \]\newline

\item \textbf{Exception :} pour un processus AR pur, les estimateurs ont une solution analytique (MCO).\newline

\item Dans le cas général (MA ou ARMA), il faut recourir à des méthodes d'\textbf{optimisation numérique}.\newline

\item Idée : calculer numériquement $\mathcal{L}(\bm{\theta})$ pour différentes valeurs de $\bm{\theta}$ et chercher la valeur qui maximise cette fonction.

\end{itemize}

\end{frame}

\begin{frame}{Recherche sur grille (\emph{grid search})}

\begin{itemize}

\item La méthode la plus simple : évaluer $\mathcal{L}(\bm{\theta})$ sur une grille de valeurs de $\bm{\theta}$.\newline

\item \textbf{Exemple :} Pour un AR(1) avec $c = 0$ et $\sigma^2 = 1$, on évalue $\mathcal{L}(\phi)$ pour $\phi \in \{-0.9, -0.8, \ldots, 0.8, 0.9\}$.\newline

\item On raffine la grille autour du maximum.\newline

\item \textbf{Avantage :} Simple et permet de visualiser la surface de vraisemblance.\newline

\item \textbf{Inconvénient :} Devient impraticable quand le nombre de paramètres augmente (malédiction de la dimension).

\end{itemize}

\end{frame}

\begin{frame}{Gradient et méthode de plus forte pente}

\begin{itemize}

\item On définit le \textbf{gradient} de la log-vraisemblance :
  \[
  \mathbf{g}(\bm{\theta}) = \frac{\partial \mathcal{L}(\bm{\theta})}{\partial \bm{\theta}}
  \]
  C'est un vecteur qui pointe dans la direction d'augmentation la plus rapide de $\mathcal{L}$.\newline

\item \textbf{Méthode de plus forte pente} (\emph{steepest ascent}) :
  \[
  \bm{\theta}^{(n+1)} = \bm{\theta}^{(n)} + \alpha_n \mathbf{g}(\bm{\theta}^{(n)})
  \]
  où $\alpha_n > 0$ est le \textbf{pas} de l'algorithme.\newline

\item On itère jusqu'à convergence : $\|\bm{\theta}^{(n+1)} - \bm{\theta}^{(n)}\| < \epsilon$ ou $\|\mathbf{g}(\bm{\theta}^{(n)})\| < \epsilon$.

\end{itemize}

\end{frame}

\begin{frame}{Calcul du gradient}

\begin{itemize}

\item Le gradient peut être calculé de deux manières.\newline

\item \textbf{Analytiquement :} On différencie $\mathcal{L}(\bm{\theta})$ par rapport à chaque élément de $\bm{\theta}$. C'est possible pour les modèles AR et MA, mais les expressions deviennent complexes pour les modèles ARMA.\newline

\item \textbf{Numériquement :} On approxime les dérivées partielles par différences finies :
  \[
  \frac{\partial \mathcal{L}}{\partial \theta_i} \approx \frac{\mathcal{L}(\bm{\theta} + h \mathbf{e}_i) - \mathcal{L}(\bm{\theta} - h \mathbf{e}_i)}{2h}
  \]
  où $\mathbf{e}_i$ est le $i$-ème vecteur de la base canonique et $h$ est un petit incrément.\newline

\item Le gradient numérique est facile à programmer mais introduit une erreur d'approximation.

\end{itemize}

\end{frame}

\begin{frame}{Maxima locaux et globaux}

\begin{itemize}

\item Si la log-vraisemblance est \textbf{unimodale}, les méthodes itératives convergent vers le maximum global, quel que soit le point de départ.\newline

\item En général, $\mathcal{L}(\bm{\theta})$ peut avoir plusieurs \textbf{maxima locaux}. L'algorithme converge alors vers le maximum local le plus proche du point de départ.\newline

\item \textbf{Stratégies pratiques :}
  \begin{itemize}
    \item Essayer plusieurs points de départ différents\newline
    \item Utiliser d'abord une recherche sur grille grossière, puis affiner avec une méthode de gradient\newline
    \item Comparer les valeurs de $\mathcal{L}(\hat{\bm{\theta}})$ obtenues à partir de différents points de départ
  \end{itemize}

\end{itemize}

\end{frame}

%==============================================================================
\section{Distribution asymptotique et tests}
%==============================================================================

\begin{frame}{La fonction de score}

\begin{itemize}

\item On appelle \textbf{fonction de score} le gradient de la log-vraisemblance :
  \[
  \mathbf{s}(\bm{\theta}) = \frac{\partial \mathcal{L}(\bm{\theta})}{\partial \bm{\theta}}
  \]
  C'est un vecteur de dimension $(p+q+2) \times 1$.\newline

\item Sous des conditions de régularité (échange dérivée/intégrale), le score a une espérance nulle évaluée en la vraie valeur :
  \[
  \E_{\bm{\theta}_0}\!\left[\mathbf{s}(\bm{\theta}_0)\right] = \mathbf{0}
  \]\newline

\item Intuition : à la vraie valeur des paramètres, la log-vraisemblance est en moyenne à son maximum.

\end{itemize}

\end{frame}

\begin{frame}{Matrice d'information de Fisher}

\begin{itemize}

\item La \textbf{matrice d'information de Fisher} est définie par :
  \[
  \mathcal{I}(\bm{\theta}) = \Var_{\bm{\theta}}\!\left[\mathbf{s}(\bm{\theta})\right] = \E_{\bm{\theta}}\!\left[\mathbf{s}(\bm{\theta})\mathbf{s}(\bm{\theta})'\right]
  \]\newline

\item Sous les conditions de régularité, on a l'\textbf{égalité de l'information} :
  \[
  \mathcal{I}(\bm{\theta}) = -\E_{\bm{\theta}}\!\left[\frac{\partial^2 \mathcal{L}(\bm{\theta})}{\partial \bm{\theta}\, \partial \bm{\theta}'}\right] = -\E_{\bm{\theta}}\!\left[\mathbf{H}(\bm{\theta})\right]
  \]
  où $\mathbf{H}(\bm{\theta})$ est la matrice \textbf{hessienne} de la log-vraisemblance.\newline

\item La matrice d'information mesure la quantité d'information que l'échantillon contient sur les paramètres.

\end{itemize}

\end{frame}

\begin{frame}{Borne de Cramér-Rao}

\begin{itemize}

\item \textbf{Borne de Cramér-Rao :} Pour tout estimateur sans biais $\hat{\bm{\theta}}$ de $\bm{\theta}$ :
  \[
  \Var(\hat{\bm{\theta}}) \geq \mathcal{I}(\bm{\theta})^{-1}
  \]
  au sens des matrices (la différence est semi-définie positive).\newline

\item L'estimateur du maximum de vraisemblance \textbf{atteint} cette borne asymptotiquement : c'est l'estimateur le plus précis (asymptotiquement) parmi les estimateurs réguliers.\newline

\item Plus la courbure de la log-vraisemblance est forte autour de $\bm{\theta}_0$ (information de Fisher élevée), plus l'estimation est précise.

\end{itemize}

\end{frame}

\begin{frame}{Distribution asymptotique de l'EMV}

\begin{itemize}

\item Sous des conditions de régularité (stationnarité, ergodicité, identifiabilité, $\bm{\theta}_0$ intérieur à l'espace des paramètres), l'EMV vérifie :\newline

\item \textbf{Convergence :}
  \[
  \hat{\bm{\theta}} \xrightarrow{p} \bm{\theta}_0 \quad \text{quand } T \to \infty
  \]\newline

\item \textbf{Normalité asymptotique :}
  \[
  \boxed{\sqrt{T}\left(\hat{\bm{\theta}} - \bm{\theta}_0\right) \xrightarrow{d} N\!\left(\mathbf{0},\; \mathcal{I}(\bm{\theta}_0)^{-1}\right)}
  \]\newline

\item En d'autres termes, pour $T$ grand :
  \[
  \hat{\bm{\theta}} \;\overset{a}{\sim}\; N\!\left(\bm{\theta}_0,\; \frac{1}{T}\mathcal{I}(\bm{\theta}_0)^{-1}\right)
  \]

\end{itemize}

\end{frame}

\begin{frame}{Erreurs standard et intervalles de confiance}

\begin{itemize}

\item En pratique, on estime la matrice de variance-covariance en remplaçant la matrice d'information par la \textbf{hessienne observée} :
  \[
  \widehat{\Var}(\hat{\bm{\theta}}) = \left[-\mathbf{H}(\hat{\bm{\theta}})\right]^{-1} = \left[-\frac{\partial^2 \mathcal{L}}{\partial \bm{\theta}\, \partial \bm{\theta}'}\bigg|_{\bm{\theta}=\hat{\bm{\theta}}}\right]^{-1}
  \]\newline

\item L'\textbf{erreur standard} du $i$-ème paramètre est :
  \[
  \text{se}(\hat{\theta}_i) = \sqrt{\left[\widehat{\Var}(\hat{\bm{\theta}})\right]_{ii}}
  \]\newline

\item Un \textbf{intervalle de confiance} à 95\% pour $\theta_i$ est :
  \[
  \hat{\theta}_i \pm 1{,}96 \times \text{se}(\hat{\theta}_i)
  \]

\end{itemize}

\end{frame}

\begin{frame}{Estimateur OPG}

\begin{itemize}

\item Une alternative à la hessienne est l'estimateur \textbf{OPG} (\emph{outer product of gradients}). On décompose la log-vraisemblance en contributions individuelles :
  \[
  \mathcal{L}(\bm{\theta}) = \sum_{t=1}^{T} \ell_t(\bm{\theta})
  \]\newline

\item La matrice d'information est estimée par :
  \[
  \widehat{\mathcal{I}}_{\text{OPG}} = \sum_{t=1}^{T} \frac{\partial \ell_t(\hat{\bm{\theta}})}{\partial \bm{\theta}} \frac{\partial \ell_t(\hat{\bm{\theta}})}{\partial \bm{\theta}'}
  \]\newline

\item Cet estimateur est facile à calculer (pas besoin de dérivées secondes), mais peut être moins précis en petit échantillon.

\end{itemize}

\end{frame}

\begin{frame}{Cas du quasi-maximum de vraisemblance}

\begin{itemize}

\item Si le modèle est mal spécifié (par exemple, on maximise une vraisemblance gaussienne alors que les erreurs ne sont pas gaussiennes), l'égalité de l'information \textbf{ne tient plus} :
  \[
  \mathbf{A} = -\E\!\left[\mathbf{H}(\bm{\theta}_0)\right] \neq \mathbf{B} = \E\!\left[\mathbf{s}(\bm{\theta}_0)\mathbf{s}(\bm{\theta}_0)'\right]
  \]\newline

\item La distribution asymptotique du QMLE est alors :
  \[
  \sqrt{T}\left(\hat{\bm{\theta}} - \bm{\theta}_0\right) \xrightarrow{d} N\!\left(\mathbf{0},\; \mathbf{A}^{-1}\mathbf{B}\,\mathbf{A}^{-1}\right)
  \]\newline

\item La matrice $\mathbf{A}^{-1}\mathbf{B}\,\mathbf{A}^{-1}$ est appelée \textbf{matrice sandwich} (ou estimateur de White/Huber). Il faut l'utiliser pour obtenir des erreurs standard robustes.

\end{itemize}

\end{frame}

\begin{frame}{Estimation de la matrice sandwich}

\begin{itemize}

\item En pratique, on estime :
  \[
  \hat{\mathbf{A}} = -\frac{1}{T}\mathbf{H}(\hat{\bm{\theta}}) = -\frac{1}{T}\frac{\partial^2 \mathcal{L}}{\partial \bm{\theta}\, \partial \bm{\theta}'}\bigg|_{\hat{\bm{\theta}}}
  \]\newline

\item Et :
  \[
  \hat{\mathbf{B}} = \frac{1}{T}\sum_{t=1}^{T} \frac{\partial \ell_t(\hat{\bm{\theta}})}{\partial \bm{\theta}} \frac{\partial \ell_t(\hat{\bm{\theta}})}{\partial \bm{\theta}'}
  \]\newline

\item L'estimateur robuste de la variance est alors :
  \[
  \widehat{\Var}_{\text{rob}}(\hat{\bm{\theta}}) = \frac{1}{T}\hat{\mathbf{A}}^{-1}\hat{\mathbf{B}}\,\hat{\mathbf{A}}^{-1}
  \]\newline

\item Si le modèle est correctement spécifié, $\hat{\mathbf{A}} \approx \hat{\mathbf{B}}$ et on retrouve l'estimateur classique.

\end{itemize}

\end{frame}

\begin{frame}{Test de Wald}

\begin{itemize}

\item On souhaite tester l'hypothèse linéaire $H_0 : \mathbf{R}\bm{\theta} = \mathbf{r}$ où $\mathbf{R}$ est une matrice $(r \times k)$ de rang $r$ et $k = p + q + 2$.\newline

\item La \textbf{statistique de Wald} est :
  \[
  W = (\mathbf{R}\hat{\bm{\theta}} - \mathbf{r})'\left[\mathbf{R}\,\widehat{\Var}(\hat{\bm{\theta}})\,\mathbf{R}'\right]^{-1}(\mathbf{R}\hat{\bm{\theta}} - \mathbf{r})
  \]\newline

\item Sous $H_0$ et asymptotiquement :
  \[
  W \xrightarrow{d} \chi^2(r)
  \]\newline

\item \textbf{Avantage :} Ne nécessite que l'estimation du modèle \textbf{non contraint}.

\end{itemize}

\end{frame}

\begin{frame}{Test de Wald : cas scalaire}

\begin{itemize}

\item Pour tester $H_0 : \theta_i = \theta_i^0$ (un seul paramètre), la statistique de Wald se simplifie en la \textbf{$t$-statistique} au carré :
  \[
  W = \left(\frac{\hat{\theta}_i - \theta_i^0}{\text{se}(\hat{\theta}_i)}\right)^2 \xrightarrow{d} \chi^2(1)
  \]\newline

\item De manière équivalente, la \emph{$t$-statistique} :
  \[
  t = \frac{\hat{\theta}_i - \theta_i^0}{\text{se}(\hat{\theta}_i)} \xrightarrow{d} N(0,1)
  \]\newline

\item \textbf{Exemple :} Pour tester si un coefficient AR ou MA est significativement différent de zéro, on calcule $t = \hat{\theta}_i / \text{se}(\hat{\theta}_i)$ et on rejette $H_0$ si $|t| > 1{,}96$ au seuil de 5\%.

\end{itemize}

\end{frame}

\begin{frame}{Test du rapport de vraisemblance}

\begin{itemize}

\item On estime le modèle sous $H_0$ (contraint, $\tilde{\bm{\theta}}$) et sous $H_1$ (non contraint, $\hat{\bm{\theta}}$).\newline

\item La \textbf{statistique du rapport de vraisemblance} (\emph{likelihood ratio}) est :
  \[
  LR = 2\left[\mathcal{L}(\hat{\bm{\theta}}) - \mathcal{L}(\tilde{\bm{\theta}})\right]
  \]\newline

\item Sous $H_0$ et asymptotiquement :
  \[
  LR \xrightarrow{d} \chi^2(r)
  \]\newline

\item \textbf{Avantage :} Ne nécessite pas le calcul de la matrice de variance-covariance.\newline

\item \textbf{Inconvénient :} Requiert l'estimation des deux modèles (contraint et non contraint).

\end{itemize}

\end{frame}

\begin{frame}{Test du rapport de vraisemblance : exemple}

\begin{itemize}

\item \textbf{Test d'un AR(1) contre un AR(2) :}
  \begin{itemize}
  \item Modèle non contraint : $Y_t = c + \phi_1 Y_{t-1} + \phi_2 Y_{t-2} + \varepsilon_t$
  \item Modèle contraint ($H_0 : \phi_2 = 0$) : $Y_t = c + \phi_1 Y_{t-1} + \varepsilon_t$
  \end{itemize}

\item On calcule :
  \[
  LR = 2\left[\mathcal{L}_{\text{AR}(2)}(\hat{\bm{\theta}}) - \mathcal{L}_{\text{AR}(1)}(\tilde{\bm{\theta}})\right]
  \]\newline

\item On rejette $H_0$ si $LR > \chi^2_{1, 0{,}95} = 3{,}84$ (seuil à 5\%, 1 degré de liberté).\newline

\item Ce test est directement applicable à la sélection de l'ordre d'un modèle ARMA.

\end{itemize}

\end{frame}

\begin{frame}{Test du multiplicateur de Lagrange}

\begin{itemize}

\item Le test du \textbf{multiplicateur de Lagrange} (ou \textbf{test du score}) utilise uniquement l'estimation sous $H_0$ :
  \[
  LM = \mathbf{s}(\tilde{\bm{\theta}})'\,\mathcal{I}(\tilde{\bm{\theta}})^{-1}\,\mathbf{s}(\tilde{\bm{\theta}})
  \]
  où $\mathbf{s}(\tilde{\bm{\theta}})$ est le score évalué à l'estimateur contraint.\newline

\item Sous $H_0$ et asymptotiquement :
  \[
  LM \xrightarrow{d} \chi^2(r)
  \]\newline

\item \textbf{Intuition :} Si $H_0$ est vraie, le score $\mathbf{s}(\tilde{\bm{\theta}})$ doit être proche de zéro. Un score élevé (en norme) fournit une preuve contre $H_0$.\newline

\item \textbf{Avantage :} Ne nécessite que l'estimation du modèle \textbf{contraint} (le plus simple).

\end{itemize}

\end{frame}

\begin{frame}{Comparaison des trois tests}

\begin{itemize}

\item Les trois tests (Wald, LR, LM) sont \textbf{asymptotiquement équivalents} sous $H_0$ : ils ont la même distribution limite $\chi^2(r)$.\newline

\item En échantillon fini, ils peuvent donner des résultats différents. On montre que :
  \[
  W \geq LR \geq LM
  \]
  Le test de Wald rejette le plus souvent, le test LM le moins.\newline

\item \textbf{Choix pratique :}
  \begin{itemize}
  \item Test de Wald : facile si le modèle non contraint est déjà estimé,
  \item Test LR : le plus couramment utilisé pour comparer des modèles emboîtés,
  \item Test LM : utile quand le modèle contraint est beaucoup plus simple.
  \end{itemize}

\end{itemize}

\end{frame}

\begin{frame}{Critères d'information}

\begin{itemize}

\item Pour comparer des modèles \textbf{non emboîtés}, on utilise des critères d'information qui pénalisent la complexité.\newline

\item \textbf{Critère d'Akaike} (AIC, 1973) :
  \[
  \text{AIC} = -2\mathcal{L}(\hat{\bm{\theta}}) + 2k
  \]
  où $k$ est le nombre de paramètres estimés.\newline

\item \textbf{Critère bayésien de Schwarz} (BIC, 1978) :
  \[
  \text{BIC} = -2\mathcal{L}(\hat{\bm{\theta}}) + k\log T
  \]\newline

\item On sélectionne le modèle qui \textbf{minimise} le critère. Le BIC pénalise davantage la complexité ($\log T > 2$ dès que $T \geq 8$) et sélectionne des modèles plus parcimonieux.

\end{itemize}

\end{frame}

\begin{frame}{Critères d'information : propriétés}

\begin{itemize}

\item \textbf{AIC :} Minimise asymptotiquement l'erreur quadratique moyenne de prévision. Tend à sélectionner des modèles légèrement sur-paramétrés.\newline

\item \textbf{BIC :} Sélectionne le vrai modèle avec probabilité tendant vers 1 quand $T \to \infty$ (\textbf{consistant}). Tend à sélectionner des modèles sous-paramétrés en petit échantillon.\newline

\item \textbf{En pratique :}
  \begin{itemize}
  \item Si l'objectif est la \textbf{prévision}, l'AIC est souvent préféré.
  \item Si l'objectif est l'\textbf{identification} du vrai modèle, le BIC est plus adapté.
  \item On estime plusieurs modèles ARMA($p$, $q$) pour $p, q \in \{0, 1, \ldots, p_{\max}\}$ et on retient celui qui minimise le critère choisi.
  \end{itemize}

\end{itemize}

\end{frame}

%==============================================================================
\section{Identification et sélection de modèles}
%==============================================================================

\begin{frame}{Approche de Box-Jenkins}

\begin{itemize}

\item Box et Jenkins (1976) proposent une procédure en quatre étapes :\newline

\item \textbf{(1) Transformation :} Transformer les données pour obtenir une série approximativement stationnaire (différenciation, logarithme).\newline

\item \textbf{(2) Identification :} Choisir les ordres $p$ et $q$ à l'aide des autocorrélations et autocorrélations partielles empiriques.\newline

\item \textbf{(3) Estimation :} Estimer les paramètres $\phi(L)$ et $\theta(L)$ par maximum de vraisemblance.\newline

\item \textbf{(4) Diagnostic :} Vérifier que le modèle estimé est compatible avec les données observées.

\end{itemize}

\end{frame}

\begin{frame}{Autocorrélations empiriques}

\begin{itemize}

\item L'\textbf{autocorrélation empirique} d'ordre $j$ est :
  \[
  \hat{\rho}_j = \frac{\hat{\gamma}_j}{\hat{\gamma}_0}
  \]
  où
  \[
  \hat{\gamma}_j = \frac{1}{T} \sum_{t=j+1}^{T} (y_t - \bar{y})(y_{t-j} - \bar{y}), \quad \bar{y} = \frac{1}{T}\sum_{t=1}^{T} y_t
  \]\newline

\item Si les données proviennent d'un processus MA($q$), alors $\rho_j = 0$ pour $j > q$. On s'attend donc à ce que $\hat{\rho}_j \approx 0$ pour $j > q$.\newline

\item Sous l'hypothèse de bruit blanc, $\hat{\rho}_j$ est approximativement distribué selon $N(0, 1/T)$. L'intervalle de confiance à 95\% est $\pm 2/\sqrt{T}$.

\end{itemize}

\end{frame}

\begin{frame}{Autocorrélations partielles empiriques}

\begin{itemize}

\item L'\textbf{autocorrélation partielle} empirique d'ordre $m$ est le dernier coefficient $\hat{\alpha}_m^{(m)}$ dans la régression :
  \[
  y_{t+1} = \hat{c} + \hat{\alpha}_1^{(m)} y_t + \hat{\alpha}_2^{(m)} y_{t-1} + \cdots + \hat{\alpha}_m^{(m)} y_{t-m+1} + \hat{e}_t
  \]\newline

\item Si les données proviennent d'un processus AR($p$), alors l'autocorrélation partielle est nulle pour $m > p$ :
  \[
  \alpha_m^{(m)} = 0 \quad \text{pour } m > p
  \]\newline

\item Sous cette hypothèse, $\Var(\hat{\alpha}_m^{(m)}) \approx 1/T$ et l'intervalle de confiance à 95\% est $\pm 2/\sqrt{T}$.

\end{itemize}

\end{frame}

\begin{frame}{Règles d'identification}

\begin{itemize}

\item \textbf{Processus MA($q$) :} Les autocorrélations $\hat{\rho}_j$ sont significativement non nulles pour $j \leq q$, puis deviennent nulles au-delà. Les autocorrélations partielles décroissent progressivement.\newline

\item \textbf{Processus AR($p$) :} Les autocorrélations partielles $\hat{\alpha}_m^{(m)}$ sont significativement non nulles pour $m \leq p$, puis deviennent nulles au-delà. Les autocorrélations décroissent progressivement (mélange d'exponentielles ou de sinusoïdes amorties).\newline

\item \textbf{Processus ARMA($p$, $q$) :} Les deux fonctions décroissent progressivement. L'identification est plus difficile et demande souvent d'essayer plusieurs spécifications.

\end{itemize}

\end{frame}

\begin{frame}{Philosophie de la parcimonie}

\begin{itemize}

\item Box et Jenkins insistent sur le principe de \textbf{parcimonie} : utiliser le moins de paramètres possible.\newline

\item Un modèle avec trop de paramètres :
  \begin{itemize}
    \item s'ajuste bien aux données historiques (sur-ajustement),
    \item mais prévoit mal hors échantillon.
  \end{itemize}

\item La découverte que des modèles ARMA avec de petites valeurs de $p$ et $q$ produisent souvent de meilleures prévisions que les grands modèles macroéconométriques a été un résultat marquant.\newline

\item En pratique, des valeurs de $p$ et $q$ inférieures ou égales à 2 ou 3 suffisent pour la plupart des applications.

\end{itemize}

\end{frame}

\begin{frame}{Résumé}

\begin{itemize}

\item Le \textbf{maximum de vraisemblance} est le principe d'estimation dominant.

\item Pour les processus \textbf{AR purs}, l'estimation se réduit aux MCO.

\item Pour les processus \textbf{MA} et \textbf{ARMA}, la vraisemblance est non linéaire et nécessite une optimisation numérique.

\item La vraisemblance \textbf{conditionnelle} simplifie les calculs mais nécessite l'inversibilité. La vraisemblance \textbf{exacte} est plus robuste.

\item L'EMV est \textbf{asymptotiquement normal} avec une variance atteignant la borne de Cramér-Rao.

\item Les tests de \textbf{Wald}, du \textbf{rapport de vraisemblance} et du \textbf{multiplicateur de Lagrange} permettent de tester des hypothèses sur les paramètres.

\item Les critères \textbf{AIC} et \textbf{BIC} complètent l'identification par les autocorrélations pour la sélection de modèles.

\end{itemize}

\end{frame}

\end{document}
